\scanpage{144}\mspage{73}

\[
 \text{(To 8)}\qquad\text{Le 2-polyèdre $P_0$ est régulier, et orientable}\tag*{(39)}
\]
\margin{Il vaut mieux séparer en deux axiomes distincts (39) -- ils ne jouent pas le $=$ rôle. De plus, il semble que dans la plupart des arguments, on n'ait besoin que du fait que $G_0^+$ soit transitif sur l'ens \unsure{$\swarrow$|.} \deleted{\unsure{$S$|$F$}} \unsure{??|?}}% slanted

ce qui fait qu'on doit pouvoir traduire notre axiomatique en termes de théorie des groupes, \unsure{d'abord|à savoir|savoir} en termes des propriétés d'un groupe $G_0^+$ qui est engendré par des générateurs $\rho_0,\rho_1,\rho_\infty$ satisfaisant
\[
 \boxed{\rho_0^n=\rho_1^2=\rho_\infty^3=1}\tag*{(40)}
\]
et où $\rho_0$ est d'ordre $n$ exactement, sous-groupe \deleted{\unsure{de|d'un}} d'indice 2 d'un groupe \inserted{$G_0$} ayant des générateurs $\tau_0,\tau_1,\tau_\infty$ reliés entre eux et aux précédents par
\[
 \left\{\begin{aligned}
 &\tau_0^2=\tau_1^2=\tau_\infty^2=1\\
 &\rho_0=\tau_\infty\tau_1,\quad \rho_1=\tau_0\tau_\infty,\quad
 \rho_\infty=\tau_1\tau_0 .
 \end{aligned}\right. \tag*{(41)}
\]

Je vais d'abord \deleted{introduire un \unsure{homo|bon} isomorphisme} \inserted{introduire le $\Lambda_n$-torseur}
\[
 \vec\Pi=\deleted{\unsure{\coprod|=}}\coprod_{i\in\Pi}\mu(i)\qquad
 \text{où }\mu(i)=\text{ens. des deux orientations de }P_i=(S,A_i,F_i)\tag*{(42)}
\]
on fait opérer dessus \deleted{$(\mathbb Z/n\mathbb Z)^*\simeq\unsure{\Lambda|\mathbb Z}$}
\[
 Z_n=\Lambda_n^{*}=(\mathbb Z/n\mathbb Z)^*\tag*{(43)}
\]
en notant que pour tout \deleted{$s\in S$ et} \inserted{$i\in\Pi$,} $\lambda\in Z'\simeq Z'_n$, \inserted{$s\in S$} on a sur $\underbrace{\operatorname{Et}_i(s)=\operatorname{Et}_{\lambda.i}(\lambda s)}_{E}$, deux structures polygônales \deleted{\unsure{$\alpha$|$a_i$}} $\alpha_i(s)$, $\alpha_{\lambda.i}(s)$, reliées par
\[
 \alpha_{\lambda.i}(s)=\alpha_i(s)^\lambda\tag*{(44)}
\]

% archive sheet number: 143
